\chapter{Foundations}
This lecture is about the security of message protocols.
There are different definitions of security, in the course of this lecture we define, what security means in this context of cryptographic protocols, and how to automatically verify the security of a protocol.

\section{Basic Definitions}
First and foremost, we define the cryptographic primitives and basic operations.

\newcommand{\encs}{\texttt{enc}^\text{s}}
\newcommand{\enca}{\texttt{enc}^\text{a}}
\newcommand{\decs}{\texttt{dec}^\text{s}}
\newcommand{\deca}{\texttt{dec}^\text{a}}
\begin{definition}[Cryptographic Primitives: Encryption]
    For \(X,Y,K \subseteq \Sigma^*\) for alphabet \(\Sigma\), we define the \emph{symmetric} encryption \(\encs : X \times K \rightarrow Y\) and decryption \(\decs : Y \times K \rightarrow X\) as some bijective functions with
    \[\decs_k \circ \encs_k = x \mapsto x.\]
    Furthermore, for \(\hat K \subseteq \Sigma^*\) and some bijective mapping to \(K\), we define the \emph{asymmetric} encryption \(\enca : X \times K \rightarrow Y\) and decryption \(\deca : Y \times \hat K \rightarrow X\) functions with property
    \[\deca_{\hat{k}} \circ \enca_{k} = x \mapsto x.\]
    In both encryption primitives, we assume the security property that there is no easily computable direct mapping \(Y \rightarrow X\), nor \(K \rightarrow \hat K\).

\end{definition}

\newcommand{\mac}{\texttt{mac}}
\newcommand{\tests}{\texttt{test}^\text{s}}
\newcommand{\testa}{\texttt{test}^\text{a}}
\newcommand{\ok}{\texttt{ok}}
\newcommand{\err}{\texttt{error}}
\newcommand{\sig}{\texttt{sig}}
\begin{definition}[Cryptographic Primitives: Signatures]
    For \(X,K,T \subseteq \Sigma^*\), we define the symmetric signature \(\mac : X \times K \rightarrow T\) and verification \(\tests : X \times K \times T \rightarrow \{\ok, \err\}\), such that
    \[\tests_k(x, \mac_k(x)) \overset{!}{=} \ok.\]
    Furthermore, for \(\hat K \subseteq \Sigma^*\) and some bijective mapping to \(K\), we define the asymmetric signature \(\sig : X \times \hat K \rightarrow T\) and verification \(\testa : X \times K \times S \rightarrow \{\ok, \err\}\), such that
    \[\testa_k(x, \sig_{\hat{k}}(x)) \overset{!}{=} \ok.\]
    Again, we assume the security property, that there is no easily computable direct mapping \(X \rightarrow T\), nor \(K \rightarrow \hat K\).

    For signatures, we often write \(\sig_k(t)\) instead of \(\sig_{\hat k}(t)\) to highlight, that the public key \(k\) is known by all, even though the private key is needed to sign \(t\).
\end{definition}

\newcommand{\hash}{\texttt{hash}}
\begin{definition}[Cryptographic Primitives: Hashing]
    For \(X, T \subseteq \Sigma^*\), we define the bijective mapping \(\hash : X \rightarrow T\), where \(\hash^{-1}\) is not easily computable.
\end{definition}

\begin{definition}[Cryptographic Primitives: Nonces and Random Numbers]
    Lastly, we assume, that we can generate perfectly random numbers or messages we call nonces, typically denoted by \(N\).
\end{definition}

\section{Formal Protocol Model}
Because we have difficulties when talking about security in an ambiguous model, we have to define a formal model to represent protocols, before defining what security even means.

In the context of this lecture, a formal model must support the following features:
1. The adversary must be able to control the routing of messages.
2. The model must be able to generate nonces.
3. The adversary may send arbitrary messages, limited only by cryptography.
4. The adversary must be able to control the scheduling.

Firstly, we have to define, how to notate messages and terms.

\newcommand{\T}{\mathcal{T}}
\newcommand{\C}{\mathcal{C}}
\newcommand{\V}{\mathcal{V}}
\newcommand{\ID}{\text{IDs}}
\begin{definition}[Terms and Messages]\label{def:terms}
    We denote the set of terms \(\T\) as the smallest set with the following properties for \(\C\) the set of constants, \(\V\) the set of variables, and \(\ID\) the set of names.
    \begin{IEEEeqnarray*}{r'c'c'c'l}
        & & \{\epsilon\} \cup \C \cup \V \cup \ID & \subseteq & \T \\
        i \in \N, a \in \ID & \Rightarrow & N_i, N_i^a, k_a, \hat k_a & \in & \T \\
        t_1, t_2 \in \T & \Rightarrow & [t_1, t_2] & \in & \T \\
        t, k \in \T & \Rightarrow & \encs_k(t) & \in & \T \\
        t \in \T \wedge a \in \ID & \Rightarrow & \enca_{k_a}(t) & \in & \T \\
        t, k \in \T & \Rightarrow & \mac_k(t) & \in & \T \\
        t \in \T \wedge a \in \ID & \Rightarrow & \sig_{k_a}(t) & \in & \T \\
        t \in \T & \Rightarrow & \hash(t) & \in & \T
    \end{IEEEeqnarray*}
    We relax this definition by allowing \(N_A\) or \(N_A''\) to be nonces generated by \(A \in \ID\) or writing \([t_0, [t_1, t_2]]\) as \(t_0, t_1, t_2\). 
    Terms can be represented as a tree or a directed acyclic graph (DAG).

    We call terms without variables a \emph{ground term} or a \emph{message}.
\end{definition}

In the tree or DAG of a term, we can trace to some position by \emph{undoing} operators specified in \cref{def:terms}, we call these terms \emph{subterms}.

\newcommand{\Sub}{\text{Sub}}
\begin{definition}[Subterms]
    For some term \(t \in \T\) and some position \(p \in \N^+\), we write \(t|_p\) as the subterm reached in \(t\) by following \(p\).

    More formally, we define the set of subterms \(\Sub(t)\) of term \(t \in \T\), as the smallest set with the following properties.
    \begin{IEEEeqnarray*}{r'c'c'c'l}
        & & t & \in & \Sub(t) \\
        t' \in \Sub(t) & \Rightarrow & \Sub(t') & \subseteq & \Sub(t) \\
        \exists t_1, t_2 : t = [t_1, t_2] & \Rightarrow & t_1, t_2 & \in & \Sub(t) \\
        \exists k, t' : t = \encs_k(t') & \Rightarrow & k, t' & \in & \Sub(t) \\
        \exists k, t' : t = \enca_k(t') & \Rightarrow & k, t' & \in & \Sub(t) \\
        \exists k, t' : t = \sig_k(t') & \Rightarrow & k, t' & \in & \Sub(t) \\
        \exists k, t' : t = \mac_k(t') & \Rightarrow & k, t' & \in & \Sub(t) \\
        \exists t' : t = \hash(t') & \Rightarrow & t' & \in & \Sub(t)
    \end{IEEEeqnarray*}
    We also define the set of subterms of a set of terms \(S \subseteq \T\) as 
    \[\Sub(S) := \bigcup_{t \in S} \Sub(t).\]
    Finally, we denote \(\Sub^1(S)\) as the set of direct subterms of \(t\).
\end{definition}

As mentioned, the adversary may generate messages only limited by cryptography.
Therefore they should not be able to send a message containing a non-derivable term.
Intuitively, we define the \emph{derivability} of some message as the \emph{Dolev-Yao closure}.

\newcommand{\DY}{\text{DY}}
\begin{definition}[Dolev-Yao Closure]
    For some \(S \subseteq \T\), we define the Dolev-Yao closure of \(S\), denoted \(\DY(S)\) as the set of messages, that are derivable from \(S\).
    The Dolev-Yao Closure of \(S\) is the smallest set with the following properties.
    \begin{IEEEeqnarray*}{r'c'r}
        S \cup \{\epsilon\} \cup \ID & \subseteq & \DY(S) \\
        {[t_1, t_2]} \in \DY(S) & \Leftrightarrow & t_1, t_2 \in \DY(S) \\
        a \in \ID, t \in \DY(S) & \Rightarrow & \enca_{k_a}(t) \in \DY(S) \\
        t, k \in \DY(S) & \Rightarrow & \encs_k(t), \mac_k(t) \in \DY(S) \\
        t, \hat k \in \DY(S) & \Rightarrow & \sig_k(t) \in \DY(S) \\
        \encs_k(t), k \in \DY(S) & \Rightarrow & t \in \DY(S) \\
        \enca_{k_a}(t), \hat k_a \in \DY(S) & \Rightarrow & t \in \DY(S) \\
        \sig_k(t) \in \DY(S) & \Rightarrow & t \in \DY(S) \\
        \mac_k(t) \in \DY(S) & \Rightarrow & t \in \DY(S) \\
        t \in \DY(S) & \Rightarrow & \hash(t)
    \end{IEEEeqnarray*}
    It should be highlighted, that this definition of the Dolev-Yao closure allows composited keys, e.g., \(\encs_{N_A,N_B}(\texttt{"Alice loves Bob"})\) to communicate a secret between Alice and Bob.
    In this example, the message is derivable from \(S \subseteq \T\), iff \(N_A, N_B \in \DY(S)\), which is the case if both nonces are derivable or if \(N_A, N_B \in S\).
\end{definition}

Now, as seen pretty quickly, the Dolev-Yao closure is always infinite, therefore we can not compute it, but we can go around this problem by handling it as a decision problem.

\begin{problem}[Derivation of a Term]\label{prob:derive}
    The DERIVE problem has the input \(S \subseteq \T, m \in \T\) and answers if \(m \in \DY(S)\) or not.
    \iffalse
    \begin{IEEEeqnarray*}{l'l}
        \textit{Problem}    & \text{DERIVE} \\
        \textit{Input}      & S \subseteq \T, m \in \T \\
        \textit{Question}   & m \in \DY(S)
    \end{IEEEeqnarray*}
    \fi
\end{problem}

It turns out, DERIVE is decidable in polynomial time using a fixed-point algorithm.
The algorithm finds a shortest derivation of \(t\) iff \(t \in \DY(S)\).
This restricts the search space by using \emph{composition} and \emph{decomposition} rules which allow to derive to and from some term.

\begin{definition}[Composition and Decomposition Rules]
    For some term \(m \in \T\), we define composition and decomposition rules as functions, that describe how \(m\) can be decomposed or composed---what to derive to with \(m\) and how to derive to \(m\).

    We denote decomposition and composition rules as \(L_d(m), L_c(m) : R \rightarrow O\), where \(R, O \subseteq \T\) are the required and obtained terms.
    This is also often written without set brackets, e.g., \(L_c([a,b]) : a,b \rightarrow [a,b]\).
    
    A whole derivation \(D(m)\) of a message \(m\) from \(S \subseteq \T\) is written as a chain of derivation rules, e.g.,
    \[S = S_0 \rightarrow_{L_0} \rightarrow S_1 \rightarrow_{L_1} \dots \rightarrow S_{n-1} \rightarrow_{L_{n-1}} \rightarrow S_n \ni m.\]
\end{definition}

The success in this notation lies in the fact, that any required non-atomic terms are always direct subterms of the term the rule is applied on/to.
More formally,
\begin{IEEEeqnarray*}{r'c'l}
    L_x(t) : R \rightarrow O & \Rightarrow & R \setminus \T_{\text{at}} \subseteq \{t\} \cup \Sub^1(t)
\end{IEEEeqnarray*}
Composition rules have exactly the terms of \(\Sub^1(m)\) as required terms.
Decomposition rules obtain only terms from \(\Sub^1(m)\).

To reduce the search space even further, we can look at the properties of some shortest derivation \(D_S(m)\) of \(m \in \T\).
\begin{IEEEeqnarray*}{r'c'l}
    L_d(t) \in D_S(m) & \Rightarrow & t \in \Sub(S) \\
    L_c(t) \in D_S(m) & \Rightarrow & t \in \Sub(S \cup \{m\})
\end{IEEEeqnarray*}
To derive \(m\) from \(S\), we only need decompositions of subterms from \(S\) or compositions of subterms of \(S\) or subterms of \(m\).

From these rules, we find the following fixed-point algorithm to solve DERIVE in polynomial time.

\pagebreak % FIXME: Remove when some text is added to the current page.

\begin{algorithm}[Algorithm for DERIVE] \hfill
\begin{lstlisting}
def derive($S$, $m$):
    $S_{\text{old}} = \emptyset$
    while $S_{\text{old}} \ne S$ do
        $S_{\text{old}} = S$
        if $\exists t \in \Sub(S \cup \{m\}) \setminus S, S \rightarrow S \cup \{t\}$ then
            $S = S \cup \{t\}$
        fi
    od
    if $m \in S$ then
        accept
    else
        reject
    fi
\end{lstlisting}
\end{algorithm}

\subsection{Formal Protocol Model: Specification}
Now we have all important steps to compose a formally sound protocol model.
We write a protocol as a set of expected and responded messages.

\begin{definition}[Receive/Send Action]
    We define a receive/send (r/s) action as some pair \((r, s) \in \T^2\) which we write \(r \rightarrow s\).
\end{definition}

Currently, we still have some problems, because a protocol should be independent of its participants and needs features not yet implemented.
For example, a protocol with some action \(\enca_{k_B}(A, x) \rightarrow \enca_{k_A}(x, N_B)\) assumes, that we have some way to handle variables, in this case \(x\), and hard codes some recognition of key and nonce origin.

To fix these issues, we introduce \emph{substitutions}.

\begin{definition}[Substitutions]
    Substitutions are functions \(\sigma : \V \rightarrow \T\) with \(\exists^{<\infty} x : \sigma(x) \ne x\).
    A substitution \(\sigma\) is called \emph{ground substitution}, if \(\forall x \in \V, \sigma(x) \ne x: \sigma(x) \text{ is message}\).

    We interpret a situation \(\sigma(x) = x\) as the case, that \(x\) is uninitialized.
    Else, we evaluate inductively for arbitrary terms\footnote{\(\T_{\text{nat}}\) is supposed to be \(\T \setminus \T_{\text{at}}\), aka., non-atomic terms.}.
    \begin{IEEEeqnarray*}{l"r'c'l}
        \text{IB: } x \in \V              & \sigma(x) & & \text{ is already defined} \\
        \text{IS: } x \in \T_{\text{nat}} \setminus \V & \sigma(f(t_1, \dots, t_n)) & := & f(\sigma(t_1), \dots, \sigma(t_n))
    \end{IEEEeqnarray*}
\end{definition}

Putting this together, we can model a protocol step, as a substitution state \(\sigma\), next r/s action, incoming message \(m\), updated substitution \(\sigma \mapsto \sigma'\), and reply message \(m'\).
